\newpage
\section{FSK-Transmitter}

% \begin{itemize}[noitemsep]
%     \item Kurze Beschreibung der Aufgabenstellung
%     \item Zugeteilter Frequenzkanal
%     \item Beschreibung der Einstellungen des CC1101 Chips mit Hilfe von Code-Snippets, zum Beispiel das Code-Snippet zur Einstellung des Kanals ist wie folgt
%     \begin{itemize}[noitemsep]
%         \item \lstinline!C1101_433.setChannel(10)  # set Channel to 10 (approx. 435 MHz)!
%     \end{itemize}
%     \item Screenshot des GNU Radio Companion Aufbaus
%     \begin{itemize}[noitemsep]
%         \item Parameter des Tiefpassfilters (warum wurden diese so gewählt?)
%     \end{itemize}
%     \item Analyse des empfangenen Signals im Zeitbereich für 2-FSK, 4-FSK und GFSK
%     \begin{itemize}[noitemsep]
%         \item Grafische Darstellung der gefilterten und nicht-gefilterten Signale
%         \item Interpretation der Ergebnisse (z.B. Unterschied FSK und GFSK)
%     \end{itemize}
%     \item Analyse des empfangenen Signals im Frequenzbereich für 2-FSK, 4-FSK und GFSK
%     \begin{itemize}[noitemsep]
%         \item Grafische Darstellung des gefilterten und nicht-gefilterten Signals
%         \item Interpretation der Ergebnisse
%     \end{itemize}
% \end{itemize}

Es soll mit dem CC1101 Transciever am eLITe-Board ein Sender aufgesetzt werden, der mit verschiedenen FSK Modulationsarten (2-FSK, 4-FSK und GFSK) Daten auf einem zugeteilten Frequenzkanal sendet. Das Signal wird dann mit dem USRP Empfangen und wie in \autoref{sec:task1} demoduliert und ausgewertet.

\subsection{Kanalzuteilung}

Der Zugeteilte Kanal ist \emph{Kanal 1}. Jeder Kanal hat einen Frequenzabstand von $\SI{200}{\kilo\hertz}$. Mittelfrequenz des Kanals beträgt daher
\begin{equation*}
    f_\text{center,1} = f_\text{0} + 200\text{kHz} \cdot \text{Channel} = 433.0\text{MHz} + 200\text{kHz} \cdot 1 = 433.2\text{MHz}
\end{equation*}

\subsection{Einstellungen des CC1101}

Nach dem Aufsetzten der eLITe Board Peripherien wird der CC1101 Transciever wie folgt konfiguriert:
\begin{lstlisting}[caption={Einstellungen des CC1101 Transcievers für Task 2}, label={lst:cc1101_task2}]
CC1101_433.setChannel(1) # Set Channel to 1 (approx. 433.2 MHz) 
CC1101_433.setBaud(1000) # Set Baudrate to 1000 bps
CC1101_433.enWhiteData(False)
CC1101_433.enFEC(False)
CC1101_433.enCRC(False)
CC1101_433.setPacketMode("PKT_LEN_FIXED")
CC1101_433.setPktLen(2)  # 2 Bytes
CC1101_433.setTXPower(7) # Max Power
CC1101_433.setSyncWord("0F0F")

MODULATIONS = ["2-FSK", "4-FSK", "GFSK"]
CC1101_433.setModulation(MODULATIONS[0])
\end{lstlisting}

Im Anschluss kann ein Datenwort \lstinline!data! definiert werden, welches dann mit \lstinline!CC1101_433.sendData(data)! in einer Schleife gesendet wird. Wurde das Datenwort am Empfänger gemessen, kann die Modulationsart mit \lstinline!CC1101_433.setModulation(MODULATIONS[i])! geändert werden, um die verschiedenen Modulationsarten zu testen.

\newpage
\subsection{Aufbau in GRC}

Der Aufbau in GRC gleicht dem aus \autoref{sec:task1}, mit dem Unterschied, dass am USRP-Source Block die Mittelfrequenz auf den zugeteilten Kanal eingestellt wurde und ein Tiefpassfilter Block hinzugefügt wurde.

\begin{figure}[h]
    \centering
    \includegraphics[width=\textwidth]{\usrpSendPath{2_grc.pdf}}
    \caption{GNU Radio Companion Aufbau für Task 2}
\end{figure}

Die USRP-Source liefert das Basisband Signal des Kanals um die Frequenz 0 Hz. Der Tiefpassfilter dient dazu, die Anteile der Nachbarkanäle zu entfernen. Die Grenzfrequenz des Filters muss dabei größer sein als die Frequenzabweichung aus \autoref{eq:f_dev} des FSK-Signals. Da die Kanalbreite $\SI{200}{\kilo\hertz}$ beträgt, wird der Tiefpassfilter mit einer Grenzfrequenz von $\SI{100}{\kilo\hertz}$ eingestellt.

\begin{figure}[h]
    \centering
    \includegraphics[width=.9\textwidth]{\usrpSendPath{2_waterfall_bpf.png}}
    \caption{Effekt des Tiefpassfilters im Frequenzbereich}
\end{figure}

\subsection{Analyse im Zeitbereich}

In \autoref{fig:task2_2fsk_gfsk} wird das demodulierte Signal von 2-FSK und GFSK gegenübergestellt. Der Verlauf des Zeitsignals ist durch den \textit{Quadrature Demod} Block direkt proportional zur Phasenänderung des empfangenen Signals. Bei 2-FSK sind steile Übergänge zwischen den Symbolen zu erkennen, da die Frequenz sprunghaft wechselt. Bei GFSK hingegen sind die Übergänge zwischen den Symbolen kontinuierlich und weisen auch kein Überschwingen auf.

\begin{figure}[h]
    \begin{minipage}[t]{.5\textwidth}
        \centering
        \includegraphics[height=4cm]{\usrpSendPath{2fsk-a5-zoom.png}}
        (a) 2-FSK Signal
    \end{minipage}\hfill
    \begin{minipage}[t]{.5\textwidth}
        \centering
        \includegraphics[height=4cm]{\usrpSendPath{gfsk-a5-zoom.png}}
        (b) GFSK Signal
    \end{minipage}
    \caption{Vergleich von 2-FSK und GFSK im Zeitbereich}\label{fig:task2_2fsk_gfsk}
\end{figure}

Bei 4-FKS sind im Zeitbereich vier verschiedene Pegel zu erkennen, welche den vier Symbolen entsprechen. Um alle Symbole zu zeigen, wurden in \autoref{fig:task2_4fsk} zwei verschiedene Datenworte \texttt{0xA5} und \texttt{0xB3} gesendet.

\begin{figure}[h]
    \begin{minipage}[t]{0.5\textwidth}
        \centering
        \includegraphics[height=4cm]{\usrpSendPath{4fsk-a5-zoom.png}}
        (a) Datenwort \texttt{0xA5}
    \end{minipage}\hfill
    \begin{minipage}[t]{.5\textwidth}
        \centering
        \includegraphics[height=4cm]{\usrpSendPath{4fsk-b3-zoom.png}}
        (b) Datenwort \texttt{0xB3}
    \end{minipage}
    \caption{4-FSK Signal im Zeitbereich}\label{fig:task2_4fsk}
\end{figure}

Welche Pegel einem Symbol angehören kann \autoref{tab:modulation_schemes} entnommen werden. Das Erhaltene Signal ist vergleichbar mit dem Beispiel aus dem Datenblatt \autoref{fig:task1_cc1101_frame}.

\begin{table}[h]
    \centering
    \caption{Symbolencoding für Verschiedene FSK-Modulationsarten \cite[S.~42, Tab.~29]{TI_CC1101_Datasheet}}
    \label{tab:modulation_schemes}
    \begin{tabularx}{\textwidth}{YYY}
        Format & Symbol & Coding \\ \midrule
        \multirow{2}{*}{2-FSK/GFSK} & \texttt{'0'} & $-f_\text{dev}$ \\
        & \texttt{'1'} & $+f_\text{dev}$ \\ \midrule
        \multirow{4}{*}{4-FSK}
        & \texttt{'01'} & $-f_\text{dev}$ \\
        & \texttt{'00'} & $-\frac{1}{3}f_\text{dev}$ \\
        & \texttt{'10'} & $+\frac{1}{3}f_\text{dev}$ \\
        & \texttt{'11'} & $+f_\text{dev}$ \\
    \end{tabularx}
\end{table}

\newpage
\subsection{Analyse im Frequenzbereich}

Für 2-FSK sind wie bisher zwei Frequenzkomponenten für die zwei Symbole zu erkennen. Diese Weisen wie in \autoref{tab:modulation_schemes} definierten Frequenzabweichungen $\pm f_\text{dev}$ vom Träger auf.

\begin{figure}[h]
    \centering
    \includegraphics[width=.6\textwidth]{\usrpSendPath{2_waterfall_2fsk.png}}
    \caption{Spektrum 2-FSK Signal}
\end{figure}

Für 4-FSK sind nun mehrere Peaks im Spektrum zu erkennen, da nun auch bis zu 4 Symbole übertragen werden, welche einem Bitmuster gemäß \autoref{tab:modulation_schemes} entsprechen. Die Frequenzkomponenten entsprechen den in \autoref{tab:modulation_schemes} definierten Abweichungen.

\begin{figure}[h]
    \centering
    \includegraphics[width=.6\textwidth]{\usrpSendPath{2_waterfall_4fsk.png}}
    \caption{4-FSK Signal}
\end{figure}

Bei GFSK gibt es ähnlich wie bei 2-FSK nur zwei Frequenzkomponenten. GSFK zeichnet sich dadurch aus, dass die Frequenz nicht sprunghaft umgeschalten wird, sondern einen kontinuierlichen übergang aufweist. Im Zeitverlauf des Spektrums in \autoref{fig:task2_gfsk} kann man erkennen, dass bei einem Symbolwechsel der Peak von $+f_\text{dev}$ nach $-f_\text{dev}$ wandert.

\begin{figure}[h]
    \centering
    \includegraphics[width=.6\textwidth]{\usrpSendPath{2_waterfall_gfsk.png}}
    \caption{GFSK Signal}\label{fig:task2_gfsk}
\end{figure}